% !TeX root = ../教材题库_全部条目_转写工作稿.tex
% Source: XJ-B1-C05-EX-P169-Y0177
% 第5章 / 5.2 任意角的三角函数 / 教材 p.169 / PDF 第176页
% OCR未校对：公式与文字均可能有误；未完成数学TeX转写；禁止直接发给学生
% 原题及相邻原解截图（核对用）：../题目原图/XJ-B1-C05-EX-P169-Y0177-01.webp
% 本文件只在转写完成后作为题干使用；OCR 条目尚不是可靠题干。
\begin{bankitem}{5.2 例题 例题6}{第5章 / 5.2 任意角的三角函数 / 教材 p.169 / PDF 第176页}{XJ-B1-C05-EX-P169-Y0177}{OCR未校对：公式与文字均可能有误；未完成数学TeX转写；禁止直接发给学生}
\begin{ocrdraft}
[教材p.169 / PDF第176页]
\par
0河已知tan仅=乐且角在第三象限，求sina，cosa\ensuremath{\cdot}
\par
解由角在第三象限知：sin<0，cos仅<0．
\par
由
\par
Slna
\par
==tana==k，得srna=kcos仅．
\par
COSa
\par
将上式代人
\par
S1n仅COS仅：一；1，
\par
得
\par
差2COS2a-+\ocrglyph{00AF}COS2
\par
1
\par
即
\par
1
\par
因此cosa=
\par
S1n．一tan仅
\par
COS仅．一
\end{ocrdraft}
\end{bankitem}
